\scanpage{128}
Or les données 1° et 2° s'expriment par des $S$-morphismes de $T$ ds des $S$-préschémas \deleted{affines} $K,L$ sur $S$ affines sur $S$. Les homomorphismes de $s_1^*(\underline L')_T$ dans $s_2^*(\underline L')_T$, ou de \deleted{$\unsure{(q_1i'')^*(\underline L')_T}$} $\mathcal O_T$ dans $(q_2i'')(\underline L')_T$, \inserted{ou de $\mathcal O_{\unsure{Z''|Z'}}$ dans $(i'\bar q_1)^*(\underline L')$,} \deleted{s'expriment également par des homomorphismes ds des \unsure{puis} $S$-schémas $M,N$ affines sur $S$, et de \unsure{la section} de $(i'\bar g')^*(\underline L')$ sur $\Sigma'$} s'expriment par des $S$-morphismes ds des $S$-préschémas $M,N,P$ \deleted{\unsure{et}} affines sur $S$, donc séparés sur $S$.

\margin{$f''$ et $Z'\longrightarrow S$ plats (i)}
\margin{$f'''$ et $Z''\text{---}S$ plats \inserted{et $\Sigma'\text{---}S$} (i)}

Les conditions écrites a) b) c) s'expriment par le fait que deux \inserted{couples} morphismes
\diagram{\begin{tikzcd}[column sep=small]
K\arrow[r,shift left=0.6ex]\arrow[r,shift right=0.6ex] & M, & L\arrow[r,shift left=0.6ex]\arrow[r,shift right=0.6ex] & N, & L\arrow[r,shift left=0.6ex]\arrow[r,shift right=0.6ex] & P
\end{tikzcd}}
s'égalisent.

Donc si \deleted{$R,S$} $U,V$ sont les noyaux des doubles flèches $K\rightrightarrows M$ et $L\rightrightarrows N\times P$ resp. [qui sont fermés ds $K,L$, donc affines sur $S$] le foncteur envisagé est représenté par $U\times V$, cqfd.
\margin{127}
